\section{Structural Decomposition}\label{sec:results}

The reduced form delivered two sharp findings---a 10--11 percent price effect that survives a battery of robustness checks, and a Gelbach residual that assigns approximately 95 percent of that effect to something other than observable entry and winner-composition channels. This section quantifies what that something is. I report three results. The intensive margin of bidding accounts for 75--85 percent of the simulated set-aside price effect across specifications. The endogenous SME entry response absorbs 60--70 percent of an otherwise larger latent shock. And the implied zero-profit entry cost per bidder is five times larger for non-SMEs than for SMEs, which rationalizes why participation responds so strongly once the protected segment opens.

\subsection{Entry response}

Table~\ref{tab:v3_entry_rates} reports the panel. Pre-policy Preg\~ao has 0.94 SMEs and 2.68 non-SMEs per auction in non-pharma, 0.55 and 2.61 in pharma. Post-policy the counts rotate: 1.87 SMEs and 1.50 non-SMEs in non-pharma, 1.22 and 1.66 in pharma. SMEs roughly double in both classes. Non-SMEs fall sharply in both classes. The set-aside therefore changes both the admissible pool and the equilibrium participation response.

\subsection{Structural decomposition of the set-aside}

Table~\ref{tab:v3_bne_decomp} gives the point estimates under the all-bidders regime and endogenous entry, which is the main specification. In non-pharma, the simulated open-regime price is $p_{S_1} = 0.759$ of the reference price. Holding the SME pool at its pre-period size and restricting the auction to SMEs raises the simulated price to $p_{S_2} = 1.114$, an intensive shift of +0.355. Letting SMEs enter at their post-policy rate pulls the simulated price down to $p_{S_3} = 1.008$, an entry-margin offset of $-0.106$. The net effect is $p_{S_3} - p_{S_1} = +0.248$.

The pharma class follows the same decomposition with larger magnitudes. $p_{S_1} = 0.655$, $p_{S_2} = 1.166$, $p_{S_3} = 0.996$, and the total effect is $+0.341$. The intensive shift is $+0.511$; the entry offset is $-0.170$. The larger magnitudes line up with the thinner baseline SME pool in pharma and the heavier auction-level heterogeneity documented in Section~\ref{sec:identification}.

The intensive share, defined as $|p_{S_2} - p_{S_1}| / (|p_{S_2} - p_{S_1}| + |p_{S_3} - p_{S_2}|)$, is 77 percent in non-pharma and 75 percent in pharma in the main specification. Table~\ref{tab:v3_sensitivity_fc} shows 74 and 81 percent under the Turnbull NPMLE. Table~\ref{tab:v3_strict_invariance} shows 85 and 79 percent under the strict-primitive-invariance benchmark. Across classes and specifications, the intensive share sits in the 75--85 percent range. The common message is that the set-aside mainly works by changing the order statistic inside the restricted auction; entry matters, but as a partial offset. Figure~\ref{fig:v3_decomposition} makes this decomposition visible.

\begin{figure}[!htbp]
\centering
\includegraphics[width=0.98\textwidth]{fig_v3_decomposition.pdf}
\caption[Bayes-Nash equilibrium price decomposition.]{\textit{Bayes-Nash equilibrium price decomposition by class.} Bars report the simulated second-order statistic $\bar{c}_{(2)}$ normalized by the item reference price, under three counterfactual scenarios. $S_1$ is the open regime at the Pre-period pool (baseline). $S_2$ imposes SME-only while holding the SME pool at its Pre-period size (intensive margin isolated). $S_3$ is the SME-only equilibrium at the observed Post-period SME pool (full policy effect with endogenous entry). Annotations report the intensive shift ($S_2-S_1$), the entry offset ($S_3-S_2$), and the total effect ($S_3-S_1$). The dotted horizontal line marks the reference price ($p^{\text{ref}} = 1$). The intensive margin dominates in both classes. Numbers from Table~\ref{tab:v3_bne_decomp}.}
\label{fig:v3_decomposition}
\end{figure}

\subsection{Entry as partial insurance}

Table~\ref{tab:v3_apv} makes the role of entry visible. V2 imposes the full set-aside while holding the SME pool fixed at its pre-period size. Relative to the realized full set-aside counterfactual V0, V2 is 161.7 percent of V0 in non-pharma and 169.7 percent in pharma. Without endogenous entry, the price effect would therefore be about 60--70 percent larger than what is observed. The same arithmetic carries over to welfare: both the allocative wedge and the MCPF term scale with the price margin, so a fixed-pool counterfactual would materially overstate the realized burden of the policy.

This is why entry is best read as the government's partial insurance. The policy creates an intensive shock by excluding non-SMEs from price formation; the induced SME entry response then dampens, but does not undo, that shock. The extensive margin is therefore not the main channel through which the set-aside operates, but it is the market's main self-correcting force once the rule is in place. Figure~\ref{fig:v3_entry_insurance} shows the size of the absorbed shock.

\begin{figure}[!htbp]
\centering
\includegraphics[width=0.98\textwidth]{fig_v3_entry_insurance.pdf}
\caption[Entry as the government's partial insurance.]{\textit{The price effect that the SME entry response absorbs.} Bars report the simulated price effect relative to the open regime, $\Delta p / p^{\text{ref}}$, under two specifications. V0 is the realized full set-aside with endogenous Post-period SME entry; it is the main-text benchmark. V2 is the counterfactual that imposes the same full set-aside but holds the SME pool at its Pre-period size. The annotation above each V2 bar reports V2's share of V0: without the induced entry response, the price effect would be 138 percent of the realized effect in non-pharma and 101 percent in pharma. The gap between the two bars is the fiscal and welfare cost that the SME entry response absorbs. Numbers from Table~\ref{tab:v3_apv}.}
\label{fig:v3_entry_insurance}
\end{figure}

\subsection{Confidence intervals and point-identification corroboration}

A cluster bootstrap at the auction level with $B=500$ replicates (Table~\ref{tab:v3_bootstrap_ci}) gives the following 95 percent CIs on $p_{S_3} - p_{S_1}$ under the all-bidders regime:
\begin{center}
\begin{tabular}{lcc}
\toprule
Class      & Mean   & 95\% CI \\
\midrule
Non-pharma & +0.236 & $[0.186, 0.289]$ \\
Pharma     & +0.305 & $[0.247, 0.364]$ \\
\bottomrule
\end{tabular}
\end{center}
The bootstrap intensive-share CI is $[64.9\%, 88.1\%]$ in non-pharma and $[62.5\%, 82.9\%]$ in pharma. Neither class looks remotely like an entry-dominated policy. The Turnbull NPMLE corroborates the main text within five percent of the total effect and keeps the same qualitative split, with the pharma intensive share if anything higher.

\subsection{Entry cost calibration}

Table~\ref{tab:v3_entry_cost} reports the zero-profit implied entry cost per bidder. In non-pharma, $\kappa^{\text{SME}} =$ R\$0.55 and $\kappa^{\neg} =$ R\$2.46. In pharma, $\kappa^{\text{SME}} =$ R\$0.11 and $\kappa^{\neg} =$ R\$0.51. The five-to-one ratio between non-SME and SME entry costs is consistent across classes and plausible given the larger documentation and compliance burden faced by bigger suppliers. These magnitudes matter because they rationalize why participation responds strongly to the protected segment even though entry is not the policy's dominant price channel.

The structural decomposition closes the Gelbach residual of Section~\ref{sec:rf}. The 95 percent of the reduced-form price effect that observable entry and winner-composition channels could not account for resolves here into the within-auction bidding response: the intensive margin operating inside a restricted pool. What the phase-2 bid-distribution moments showed directly---bidders responding to a restricted pool by crowding toward the reference---is now quantified structurally at 75--85 percent of the effect.
